\documentclass[letterpaper, 12pt]{article}

\title{PSU CS162: Syntax Basics \color{red}\textbf{[KEY]}\normalcolor}
\author{Amiel L. Iliesi}
\date{2026-03-27}

\usepackage{tabularx}
\usepackage[hidelinks]{hyperref}
\usepackage{minted}
\usepackage[x11names]{xcolor}
\usepackage{amssymb}
\usepackage{ulem}

% commands
\newcommand{\question}[2][2em]{\textbf{#2}\vspace{#1}\par}
\newcommand{\ans}[1]{\color{Green4}\textbf{#1}\normalcolor}

% environments
\newenvironment{ansblock}{\color{Green4}}{\normalcolor\vspace{1em}}

% global options
\setlength{\parindent}{0pt}
\usemintedstyle{monokai}

\begin{document}

\maketitle

\tableofcontents

\pagebreak

\section{Python Basics}

\textbf{\emph{Write Python code that does the following...}}

\vspace{2em}

\question[0pt]{Initialize a variable with your name, and print out
    \texttt{"Hello, <your name>"}.}

\begin{minted}[bgcolor=darkgray]{python}
name = "Amiel"
print("Hello, " + name)
\end{minted}

\question[0pt]{Set up a main call and print \texttt{"Hello, World!"}.}

\begin{minted}[bgcolor=darkgray]{python}
if __name__ == "__main__":
    print("Hello, World!")
\end{minted}

\question[0pt]{Save a string of your choice to a variable. Write a conditional statement
    that prints if the saved string is either \texttt{"String A"},
    \texttt{"String B"}, or none of the above.}

\begin{minted}[bgcolor=darkgray]{python}
my_string = "Whatever I want!"

if my_string == "String A":
    print(my_string + " matches \"String A\"!")
elif my_string == "String B":
    print(my_string + " matches \"String B\"!")
else:
    print(my_string + " doesn't match any of the above...")
\end{minted}


\begin{minipage}{\textwidth}
\question{Connect the following operations to their definition.}
\begin{tabularx}{0.95\linewidth}{| X | X |}
    \hline
    \textbf{a.} \texttt{*} & exponentiation \hfill \ans{f.}\\
    \hline
    \textbf{b.} \texttt{/} & modulo assignment \hfill \ans{d.}\\
    \hline
    \textbf{c.} \texttt{+=} & addition assignment \hfill \ans{c.} \\
    \hline
    \textbf{d.} \texttt{\%=} & integer division \hfill \ans{g.} \\
    \hline
    \textbf{e.} \texttt{\^} & float division \hfill \ans{b.} \\
    \hline
    \textbf{f.} \texttt{**} & bitwise exclusive-or \hfill \ans{e.} \\
    \hline
    \textbf{g.} \texttt{//} & multiplication \hfill \ans{a.} \\
    \hline
\end{tabularx}
\end{minipage}

\pagebreak

\section{C++ Basics}

\question[0em]{Write the setup for the \texttt{main} function.}

\begin{minted}[bgcolor=darkgray]{c++}
int main()
{
    return 0;
}
\end{minted}

\question[0em]{Write a function that returns the square of the input.}

\begin{minted}[bgcolor=darkgray]{c++}
int square(int x)
{
    return x * x;
}
\end{minted}

\question[0em]{Output \texttt{"Hello, World!"}.}

\begin{minted}[bgcolor=darkgray]{c++}
std::cout << "Hello, World!" << std::endl;
\end{minted}

\question[0em]{Explain the difference between a function implementation and a
    function prototype.}
\begin{ansblock}
A \uline{function implementation} describes \textbf{how} a function works,
and a \uline{function prototype} describes the \textbf{call signature}.
\end{ansblock}

\question[0em]{What is a function signature?}
\begin{ansblock}
A \uline{function signature} are the parts that \textbf{uniquely identify}
one function from another. In C++ the signature is the combination of the
\textbf{function name} with the
\textbf{function parameters: number of, types, and order}.
\end{ansblock}

\question[0em]{Write the bash command to compile \texttt{main.cpp} into \texttt{main.out}}
\begin{minted}[bgcolor=lightgray]{bash}
g++ main.cpp -o main.out
\end{minted}

\section{C++ VS. Python}

\begin{tabular}{| l | c | c | c |}
    \hline
    \textbf{feature} & \textbf{Python?} & \textbf{C++?} \\
    \hline
    Code scope is indent sensitive & \ans{\checkmark} & \\
    \hline
    Code scope is braced & & \ans{\checkmark} \\
    \hline
    Strings can use either single or double quote & \ans{\checkmark} & \\
    \hline
    Strings are double quote, chars are single quote & & \ans{\checkmark} \\
    \hline
    Loosely typed & \ans{\checkmark} & \\
    \hline
    Strongly typed & & \ans{\checkmark} \\
    \hline
    Interpreted & \ans{\checkmark} & \\
    \hline
    Compiled & & \ans{\checkmark} \\
    \hline
\end{tabular}

\end{document}